% GENERATED FILE. Edit canonical lessons, then run npm run generate:books.
% canonical-source-hash: fnv1a64:d0b9304655291a24
% canonical-lessons: GE-C323-schaelen, GE-C323-hacken, GE-C323-wuerfeln, GE-C323-raspeln, GE-C323-halbieren

\chapter{Schälen und schneiden}
\label{ch:ge-schalen-und-schneiden}
\hlchaptermodality{323}

\begin{chapteropening}
\textbf{By the end of this chapter:} \emph{I can describe getting vegetables ready: to peel, to chop, to dice, to grate and to halve.}

Complete the last lesson of chapter 323: I can describe getting vegetables ready: to peel, to chop, to dice, to grate and to halve.
\end{chapteropening}

\section[schälen]{schälen — to peel}

\label{lesson:GE-C323-schaelen}



\begin{quote}
Before the new one: say the German for to measure out, then the German for the bowl.
\end{quote}

\subsection*{You'll want to know: schälen}
\textbf{schälen} — \textquotedblleft{}to peel\textquotedblright{}.

\begin{quote}
Zuerst schälte er die Kartoffeln.
\end{quote}

\begin{quote}
\emph{First he peeled the potatoes.}
\end{quote}

Built on \emph{die Schale} — peel, skin; also a bowl.

\begin{cousinweb}
\emph{schälen} is built on \emph{Schale} \textquotedblleft{}peel, shell\textquotedblright{}, an inherited Germanic word from Old High German \emph{scāla} \textquotedblleft{}shell, bowl\textquotedblright{} (EtymDB).
\end{cousinweb}

\begin{grammarlens}[title={What you've built}]
One more word for telling, explaining and describing at B1.
\end{grammarlens}

\subsection*{Your turn}
\begin{itemize}
\raggedright
  \item \emph{Say it:} \emph{schälen}
  \item \emph{Say it:} \emph{schälen}, once more
  \item \emph{Say it:} a sentence of your own with \emph{schälen}
  \item \emph{Recall:} say the German for to measure out, then the German for the bowl, then say \emph{schälen} again
\end{itemize}

\begin{tcolorbox}[breakable,colback=teal!4,colframe=teal!35!black,title={Before you move on}]
What does schälen mean? (\textquotedblleft{}To peel\textquotedblright{}.) Say it once more.
\end{tcolorbox}
\section[hacken]{hacken — to chop}

\label{lesson:GE-C323-hacken}



\begin{quote}
Before the new one: say the German for the bowl, then the German for to peel.
\end{quote}

\subsection*{You'll want to know: hacken}
\textbf{hacken} — \textquotedblleft{}to chop\textquotedblright{}.

\begin{quote}
Sie hackte die Zwiebeln ganz klein.
\end{quote}

\begin{quote}
\emph{She chopped the onions very fine.}
\end{quote}

Also to chop wood: \emph{Holz hacken}.

\begin{cousinweb}
\emph{hacken} is usually traced to an inherited Germanic verb, a cousin of English \emph{hack}; EtymDB lists no parent.
\end{cousinweb}

\begin{grammarlens}[title={What you've built}]
One more word for telling, explaining and describing at B1.
\end{grammarlens}

\subsection*{Your turn}
\begin{itemize}
\raggedright
  \item \emph{Say it:} \emph{hacken}
  \item \emph{Say it:} \emph{hacken}, once more
  \item \emph{Say it:} a sentence of your own with \emph{hacken}
  \item \emph{Recall:} say the German for the bowl, then the German for to peel, then say \emph{hacken} again
\end{itemize}

\begin{tcolorbox}[breakable,colback=teal!4,colframe=teal!35!black,title={Before you move on}]
What does hacken mean? (\textquotedblleft{}To chop\textquotedblright{}.) Say it once more.
\end{tcolorbox}
\section[würfeln]{würfeln — to dice; to throw dice}

\label{lesson:GE-C323-wuerfeln}



\begin{quote}
Before the new one: say the German for to peel, then the German for to chop.
\end{quote}

\subsection*{You'll want to know: würfeln}
\textbf{würfeln} — \textquotedblleft{}to dice; to throw dice\textquotedblright{}.

\begin{quote}
Dann würfelte er die Karotten.
\end{quote}

\begin{quote}
\emph{Then he diced the carrots.}
\end{quote}

To cut into small cubes. With dice, \emph{würfeln} is to throw them.

\begin{cousinweb}
\emph{würfeln} is built on \emph{Würfel} \textquotedblleft{}a dice\textquotedblright{}, Old High German \emph{wurfil}, Proto-Germanic \emph{wurpilaz} (EtymDB), from the root of \emph{werfen} \textquotedblleft{}to throw\textquotedblright{}.
\end{cousinweb}

\begin{grammarlens}[title={What you've built}]
One more word for telling, explaining and describing at B1.
\end{grammarlens}

\subsection*{Your turn}
\begin{itemize}
\raggedright
  \item \emph{Say it:} \emph{würfeln}
  \item \emph{Say it:} \emph{würfeln}, once more
  \item \emph{Say it:} a sentence of your own with \emph{würfeln}
  \item \emph{Recall:} say the German for to peel, then the German for to chop, then say \emph{würfeln} again
\end{itemize}

\begin{tcolorbox}[breakable,colback=teal!4,colframe=teal!35!black,title={Before you move on}]
What does würfeln mean? (\textquotedblleft{}To dice; to throw dice\textquotedblright{}.) Say it once more.
\end{tcolorbox}
\section[raspeln]{raspeln — to grate (coarsely)}

\label{lesson:GE-C323-raspeln}



\begin{quote}
Before the new one: say the German for to chop, then the German for to dice.
\end{quote}

\subsection*{You'll want to know: raspeln}
\textbf{raspeln} — \textquotedblleft{}to grate (coarsely)\textquotedblright{}.

\begin{quote}
Sie raspelte den Käse für die Pizza.
\end{quote}

\begin{quote}
\emph{She grated the cheese for the pizza.}
\end{quote}

\emph{Reiben} is finer grating.

\begin{cousinweb}
\emph{raspeln} is usually traced to Middle High German \emph{raspen} \textquotedblleft{}to scrape together\textquotedblright{}; EtymDB lists no parent.
\end{cousinweb}

\begin{grammarlens}[title={What you've built}]
One more word for telling, explaining and describing at B1.
\end{grammarlens}

\subsection*{Your turn}
\begin{itemize}
\raggedright
  \item \emph{Say it:} \emph{raspeln}
  \item \emph{Say it:} \emph{raspeln}, once more
  \item \emph{Say it:} a sentence of your own with \emph{raspeln}
  \item \emph{Recall:} say the German for to chop, then the German for to dice, then say \emph{raspeln} again
\end{itemize}

\begin{tcolorbox}[breakable,colback=teal!4,colframe=teal!35!black,title={Before you move on}]
What does raspeln mean? (\textquotedblleft{}To grate (coarsely)\textquotedblright{}.) Say it once more.
\end{tcolorbox}
\section[halbieren]{halbieren — to halve, to cut in half}

\label{lesson:GE-C323-halbieren}



\begin{quote}
Before the new one: say the German for to dice, then the German for to grate.
\end{quote}

\subsection*{You'll want to know: halbieren}
\textbf{halbieren} — \textquotedblleft{}to halve, to cut in half\textquotedblright{}.

\begin{quote}
Er halbierte die Tomaten.
\end{quote}

\begin{quote}
\emph{He cut the tomatoes in half.}
\end{quote}

Built on \emph{halb} with the borrowed ending \emph{-ieren}: \emph{hat halbiert}.

\begin{cousinweb}
\emph{halbieren} is built on \emph{halb} \textquotedblleft{}half\textquotedblright{}, Old High German \emph{halb}, Proto-Germanic \emph{halbaz} (EtymDB), with the ending \emph{-ieren}.
\end{cousinweb}

\begin{grammarlens}[title={What you've built}]
One more word for telling, explaining and describing at B1.
\end{grammarlens}

\subsection*{Your turn}
\begin{itemize}
\raggedright
  \item \emph{Say it:} \emph{halbieren}
  \item \emph{Say it:} \emph{halbieren}, once more
  \item \emph{Say it:} a sentence of your own with \emph{halbieren}
  \item \emph{Recall:} say the German for to dice, then the German for to grate, then say \emph{halbieren} again
\end{itemize}

\begin{tcolorbox}[breakable,colback=teal!4,colframe=teal!35!black,title={Before you move on}]
What does halbieren mean? (\textquotedblleft{}To halve, to cut in half\textquotedblright{}.) Say it once more.
\end{tcolorbox}
